\subsection{Ambiente}
Per garantire la massima naturalezza durante l'interazione e simulare un contesto di utilizzo reale, si è deciso di condurre gli esperimenti in luoghi familiari e confortevoli per gli utenti, mettendoli così a proprio agio. È stata prestata particolare attenzione affinché gli ambienti scelti fossero privi di distrazioni o interferenze esterne, permettendo ai partecipanti di concentrarsi esclusivamente sullo svolgimento dei compiti assegnati.

Al fine di mantenere le condizioni di test rigorosamente identiche per tutti i soggetti, le prove sono state eseguite sul computer portatile dello sperimentatore, dotato di mouse esterno per uniformare e agevolare l'interazione. Questa scelta ha consentito di standardizzare l'ambiente software (utilizzando il medesimo browser, \textit{Google Chrome}) e ha facilitato l'esecuzione in background degli strumenti necessari per il monitoraggio e la misurazione delle interazioni (come le estensioni per il tracciamento dei click e dei tempi).

Per quanto riguarda la somministrazione dei questionari, ci si è avvalsi della piattaforma Google Moduli. I form sono stati compilati in formato digitale, una soluzione che ha permesso di evitare l'uso di supporti cartacei, velocizzando notevolmente la fase di compilazione e snellendo la successiva elaborazione e analisi dei dati. 

Infine, per assicurare che la preparazione della postazione fosse uniforme per ogni partecipante, gli sperimentatori hanno seguito scrupolosamente una specifica \textit{checklist} (Tabella \ref{checklist}) preliminare prima di dare inizio a ciascuna sessione.